% !TeX root = ../main.tex
\chapter*{How to Use This Book}
\addcontentsline{toc}{chapter}{How to Use This Book}
\markboth{How to Use This Book}{How to Use This Book}
\label{ch:how-to-use}

\noindent\textbf{Chapter structure.}\enspace Each chapter opens by stating the
problem it addresses and where it sits in the design chain, develops the
method, and closes by identifying what the method cannot establish. Worked
examples are integrated into the argument rather than segregated, because the
point of an example here is usually to expose an assumption rather than to
demonstrate arithmetic.

\noindent\textbf{Presentation.}\enspace The book uses no framed boxes and no
colour. Cautions, worked calculations and points of emphasis run in the main
text, introduced where necessary by a bold lead-in. Anything important enough
to box is important enough to belong in the argument.

\noindent\textbf{Three reading paths.}\enspace The book can be read straight
through, but three shorter routes serve particular needs.

A \emph{characterisation} path runs Part~\ref{part:characterisation} in full,
then Chapters~\ref{ch:rock-mechanics} and~\ref{ch:rock-mass-parameters}. It
answers the question of how observations become design parameters, and it is
the route for an engineer preparing or reviewing a factual and interpretative
ground investigation report.

A \emph{design} path begins with Chapter~\ref{ch:rock-mass-parameters} for the
parameters, takes Chapter~\ref{ch:ground-behaviour} for the behaviour type,
then runs Part~\ref{part:design} in order. It answers the question of which
method should govern a given support decision, and it is the route for
reference and detailed design.

A \emph{delivery} path takes Chapter~\ref{ch:challenges} and
Chapter~\ref{ch:digital-twin} for the framing, then
Part~\ref{part:construction} and Part~\ref{part:performance}. It answers the
question of how a design survives contact with construction, and it is the
route for construction-stage engineers and asset owners.

\noindent\textbf{Cross-references.}\enspace Every method is referred back to
the chapter that establishes the quantities it consumes, so that a parameter
can always be traced to its origin. Calculation procedures that can be
followed by hand are collected in Appendix~\ref{app:calculations}; closed-form
solutions and design charts in Appendix~\ref{app:closed-form}; and typical
parameter ranges, for checking rather than for designing, in
Appendix~\ref{app:constants}.

\noindent\textbf{Software.}\enspace No method in this book requires a
particular program. Every calculation is stated so that it can be followed by
hand, and Appendix~\ref{app:calculations} sets out the longer ones as numbered
sequences for that purpose. Where a calculation is tedious by hand or
instructive when varied, the text points to \companionname, described in
Appendix~\ref{app:software} and released openly at \companionlink. It is
provided so that the worked examples can be reproduced and varied, not as a
design tool: it contains no judgement, and it returns an answer whether or not
its assumptions apply.
